\section{Introduction}
\label{sec:intro}

The Districting Integrity Act (DIA) specifies a single deterministic
random seed for the METIS partitioner, derived from the census release
identifier via SHA-256: $s_0 = \mathrm{SHA256}(\mathtt{census\_release\_id}
\,||\, \mathtt{DIA\_SEED\_V1})$.
This design choice raises a fundamental question: \emph{does the specific
seed matter?}

METIS is a heuristic algorithm; different seeds produce different
partitions via their effect on the initial coarsening and FM refinement
steps~\citep{karypis1998}.  If seed sensitivity is high --- meaning small
changes in the seed produce large changes in the partition quality or
partisan outcomes --- then a single-seed specification would be legally
problematic: the party that controls the seed (through control of census
timing or versioning) would hold an outsized influence over the map.

Conversely, if seed sensitivity is low --- meaning the solution space
is well-connected and the DIA seed lands near the global minimum ---
then the single-seed specification is defensible on two grounds:
(1) the chosen map is near-optimal in compactness, and (2) alternative
seeds produce similar partisan outcomes.

\subsection*{Contributions}

This paper makes four contributions:

\begin{enumerate}
  \item We characterise the solution space of METIS recursive bisection
    by sweeping 10,000 seeds per state across all 50 states (500 thousand
    total METIS calls), measuring normalised edge cut and partisan outcomes
    for each seed.
  \item We compute the coefficient of variation (CV) of the normalised
    edge cut across seeds by state category (single-district, small, medium,
    large), finding CV $< 2\%$ for 48 of 50 states.
  \item We establish that the DIA seed is within $2.9\% \pm 1.7\%$ of the
    minimum-edge-cut seed across all states, statistically indistinguishable
    from a randomly chosen seed.
  \item We measure partisan variance across seeds for Wisconsin ($k{=}8$)
    and North Carolina ($k{=}14$), the two highest-variance states, finding
    that seed selection drives less partisan variance than algorithm selection
    (B.0: 12.8pp gap between algorithm families).
\end{enumerate}

\subsection*{Connection to ConvergenceSweep}

These results motivate the ConvergenceSweep procedure (B.16), which
certifies convergence to a local minimum by running additional seeds
until $T = 600$ consecutive seeds produce no improvement.
The empirical data in this paper --- specifically, the per-state
empirical CDF of the last-improvement seed index --- justifies the
$T = 600$ threshold: 47 of 50 states show their last improvement
before seed index 600, and the three exceptions (GA, NC, WI) are
handled by the extended $T = 1{,}200$ sweep.

\subsection*{Paper structure}

Section~\ref{sec:related} reviews prior work on solution space
characterisation for redistricting.
Section~\ref{sec:methodology} describes the sweep design and metrics.
Section~\ref{sec:results} presents results by state category and
partisan implications.
Section~\ref{sec:discussion} interprets the findings for the DIA.
Section~\ref{sec:conclusion} concludes.
